\section{Growing finite memory and finite-fit non-evidence}
\label{sec:growing}

\Cref{cor:noprimecoefficient} freezes the dimension before \(k\) tends to
infinity.  Finite experiments reverse those quantifiers unless the model is
held fixed.  The reversal admits an exact universal construction.

\begin{theorem}[Bilinear nilpotent memorizer]
\label{thm:bilmemorizer}
Fix \(N\ge1\) and arbitrary
\(\eta_1,\ldots,\eta_N\in\mathbb F\).  Let
\(J_N\in M_N(\mathbb F)\) satisfy
\[
 J_Ne_j=e_{j+1}\quad(j<N),\qquad J_Ne_N=0.
\]
With
\[
        v=e_1,\qquad B=I,\qquad
        u=\sum_{j=1}^N\eta_je_j,
\]
one has
\[
 u^{\mathsf T}J_N^{k-1}Bv
 =
 \begin{cases}
 \eta_k,&1\le k\le N,\\
 0,&k>N.
 \end{cases}
\]
\end{theorem}

\begin{proof}
For \(k\le N\), \(J_N^{k-1}e_1=e_k\), so the contraction reads the \(k\)-th
coordinate of \(u\).  For \(k>N\), \(J_N^{k-1}=0\).
\end{proof}

\begin{theorem}[Trace nilpotent memorizer]
\label{thm:trmemorizer}
With the same \(J_N\), let \(B_\eta\) have entries
\[
        (B_\eta)_{1k}=\eta_k,\qquad 1\le k\le N,
\]
and all other entries zero.  Then
\[
 \operatorname{tr}(J_N^{k-1}B_\eta)
 =
 \begin{cases}
 \eta_k,&1\le k\le N,\\
 0,&k>N.
 \end{cases}
\]
\end{theorem}

\begin{proof}
The only nonzero contribution to the trace is
\[
        (J_N^{k-1})_{k1}(B_\eta)_{1k}=\eta_k.
\]
Nilpotence gives zero after \(N\).
\end{proof}

The two constructions expose where the apparent information resides:
\(\eta\) is stored directly in the endpoint covector \(u\) or the first row
of \(B_\eta\).  Nothing in the successor--divisor source selects those
entries.

\begin{corollary}[Arbitrary-prefix universality]
\label{cor:provestoomuch}
For every cutoff \(N\), the same nilpotent architecture realizes each of:
\[
\begin{array}{c}
\text{prime indicator},\quad
\text{square indicator},\quad
\text{powers of two},\quad
\text{Fibonacci membership},\\
\text{a seeded bit string},\quad
\text{a hash-derived string},\quad
\text{arbitrary signed rational values}.
\end{array}
\]
Therefore finite-prefix success is
\[
        \text{\normalfont\scshape Control}
        \mid
        \text{\normalfont\scshape Proves Too Much}.
\]
\end{corollary}

\begin{table}[t]
\centering
\caption{The fixed and growing statements have different quantifier order.
Only the first tests an infinite source-derived response.}
\label{tab:quantifiers}
\small
\begin{tabularx}{\textwidth}{L{3.1cm}YY}
\toprule
Regime & Quantifiers & Interpretation \\
\midrule
fixed finite dimension
& \(\forall d,A,B,u,v\), no infinite prime-only exact support
& asymptotic no-go from SML \\
growing cutoff model
& \(\forall N,\eta_{1:N}\), an \(N\)-dimensional realization exists
& target vector is memorized \\
valid finite audit
& freeze \(d,A,B,u,v\), then enlarge evaluation cutoff
& tests the fixed theorem without changing the model \\
\bottomrule
\end{tabularx}
\end{table}

\begin{remark}[Why a large prime prefix is not special]
The memorizer fits every bit pattern with zero residual.  Increasing \(N\)
and rebuilding \(J_N,u,B_\eta\) merely enlarges the stored table.  A finite
prime experiment becomes informative only if the dimension and all
parameters are frozen first and evaluated beyond their construction range.
Matched arbitrary-target controls are mandatory even then.
\end{remark}

This control also guards the later countable construction.  Infinite memory
can encode a total decider; the question is not whether primality is
computable, but whether the computation is source-derived and visible to
the same recurrent determinant.
